\begin{proposition}[The quadratic free directions along $\Sigma$ and $\Sigma^*$]
\label{prop:c11-freexi}
\begin{itemize}
\item (i) The quadratic local functionals of degree $-1$, of Weyl weight $\tfrac92$ in the fields and derivatives, linear in $\xi\in S_+\otimes W$, translation invariant and covariant under $\mathrm{Spin}(4)\times\mathrm{Spin}(6)\times\mathfrak g$, are spanned by $T_1$, $T_2$, $T_3$, $T_4$ and $\int\big((d^*A|_V)\xi,\lambda^*\big)$; their $Q$-closed combinations are spanned by $S^V_2(\xi)$ and $S^\perp_2(\xi)$, and the $Q$-exact ones are $0$.
\item (ii) The arity-one equation along $\xi$, the arity-one part along $\xi$ of $s(\hbar F_1)=\hbar\,\mathrm{Ob}$ (Theorem \ref{thm:c8-obstruction} (ii)), is
\begin{equation}
(Q+\{I,-\})\,\hbar F_1^{(1)}(\xi)+\big\{S^{(1)}(\xi),\hbar F_1^{(0)}\big\}=\hbar\,\mathrm{Ob}^{(1)}(\xi).
\label{eq:c11-arity1xi}
\end{equation}
At fixed $\hbar F_1^{(0)}$ two solutions differ, in the quadratic sector, by $\theta^\xi_V$ units times $S^V_2(\xi)$ plus $\theta^\xi_\perp$ units times $S^\perp_2(\xi)$, and the cubic part of \eqref{eq:c11-arity1xi} then contains $\theta^\xi_V\{I_3,S^V_2(\xi)\}+\theta^\xi_\perp\{I_3,S^\perp_2(\xi)\}$ units, $I_3$ the cubic part of $I$. For $\theta^\xi_V=\theta^\xi_\perp=:\theta^\xi$ the direction extends to $\theta^\xi$ units times $S^{(1)}(\xi)$, the arity-one value of $G_B$ (Definition \ref{def:c8-B}) and of the field redefinition by the R-rotation $S^{(1)}(R_\xi)$, $[R_\xi,\xi]=\xi$ (Theorem \ref{thm:c8-twist}). The cubic part fixes $\theta^\xi_V-\theta^\xi_\perp$ when the class of $\{I_3,S^V_2(\xi)\}$ modulo the images under $Q$ of the cubic local functionals of degree $-1$ is nonzero, as for the sector coefficients along $\Delta$ (Proposition \ref{prop:c11-arity1} (iii)).
\item (iii) The one-body kernel of a letter $X$ along $\xi$ is $-\{\hbar F_1^{(1)}(\xi),X\}$; the free directions add $-\theta^\xi_V$ units times the arrows of $S^V_2(\xi)$ and $-\theta^\xi_\perp$ units times those of $S^\perp_2(\xi)$, which are $\{S^V_2(\xi),A|_V\}=\Gamma(\xi,\lambda)|_V$ and $\{S^\perp_2(\xi),A|_{W_{\mathbb R}^\perp}\}=\Gamma(\xi,\lambda)|_{W_{\mathbb R}^\perp}$ on the gauge field and the scalars and $\tfrac12\rho(dA|_V)\xi$ and $\tfrac12\rho(dA|_{W_{\mathbb R}^\perp})\xi$ on $\lambda$ (Lemma \ref{lem:c11-lines} (i)).
\item (iv) Along $\Sigma^*$, with the position included at weight $-1$ and $\eta$ at weight $\tfrac12$, the $Q$-closed quadratic directions of degree $-1$ and Weyl weight $4$, linear in $\eta$, are spanned by $\sum_\mu\int x^\mu S^V_2(\rho(e_\mu)\eta)$ and $\sum_\mu\int x^\mu S^\perp_2(\rho(e_\mu)\eta)+\int s_0$, whose sum is the quadratic part of $S^{(1)}(\eta)$.
\end{itemize}
The same holds along $S_-\otimes W^*$ with the chiralities exchanged.
\end{proposition}

\begin{proof}
(i) The field contents of degree $-1$ and field weight $\tfrac92$ linear in a constant spinor are $(\lambda,A^*)$ without derivatives and $(A,\lambda^*)$ with one derivative, since $c$ and $c^*$ would need a further weight $\tfrac12$; the covariant pairings are $\Gamma(\xi,\lambda)$ with $A^*$ in its parts along $V$ and $W_{\mathbb R}^\perp$, and $\rho(\partial A)\xi$ with $\lambda^*$, $\partial A$ being $dA|_V$, $d^*A|_V$ or the mixed part $dA|_{W_{\mathbb R}^\perp}$, the volume element acting by a sign on each chirality. On functionals $Q$ substitutes $A\mapsto-dc$, $A^*\mapsto-\tfrac12d^*dA$, $\lambda^*\mapsto{\not d}\lambda$, $c^*\mapsto d^*(A^*|_V)$ and annihilates $c$ and $\lambda$ (Proposition \ref{prop:c1-Q}). $Q\int((d^*A|_V)\xi,\lambda^*)$ contains $-\int((d^*dc)\xi,\lambda^*)$, produced by no other structure, so its coefficient vanishes in a closed combination; $Q$ preserves the split into the parts along $V$ and along $W_{\mathbb R}^\perp$, and $S^V_2(\xi)$, $S^\perp_2(\xi)$ are closed, being the two parts of the quadratic part of $\{S,S^{(1)}(\xi)\}=0$. A $Q$-exact functional is $Q$ of a functional of degree $-2$ and field weight $\tfrac92$, and two antifields have weight at least $5$, while $c^*$ needs a further field of weight $\tfrac12$. Both evaluations of Remark \ref{rem:c11-oddmethod} give the same ranks on plane-wave normal forms. (ii) Theorem \ref{thm:c8-obstruction} (iii) restricted to arity one along $\xi$ at fixed arity-zero part. The cubic part of $s\big(\theta^\xi_VS^V_2+\theta^\xi_\perp S^\perp_2\big)$ is the bracket with $I_3$, and $\{I_3,S^V_2(\xi)+S^\perp_2(\xi)\}=-QS_3(\xi)$ with $S_3(\xi)$ the cubic part of $S^{(1)}(\xi)$, the cubic part of $\{S,S^{(1)}(\xi)\}=0$; $G_B^{(1)}(\xi)=S^{(1)}(B\xi)=S^{(1)}(\xi)$ for $\xi\in S_+\otimes W$. (iii) Theorem \ref{thm:c11-classes} (iii) and Lemma \ref{lem:c11-lines} (i). (iv) The same enumeration with $\int x^\mu$ times the structures of (i) at $\rho(e_\mu)\eta$, $\int s_0$ and $\int(\rho(A|_V)\eta,\lambda^*)$, and the exact ranks of $Q$ on plane-wave normal forms; $s_0$ enters the closed direction together with the part along $W_{\mathbb R}^\perp$.
\end{proof}

\begin{proposition}[Renormalisation freedom of $Q_1$, $S_1$ and $K_1$]
\label{prop:c11-freedom}
Under the changes of Proposition \ref{prop:c8-laws} the components change as follows; $\{S^{(1)}(x),-\}$ is the classical transformation along $x$.
\begin{itemize}
\item (i) A field redefinition $F_{\mathfrak{psl}(4|4)}\mapsto F_{\mathfrak{psl}(4|4)}+\hbar s\phi$ changes the component along $x$ by $-\{(s\phi)^{(1)}(x),-\}$, a derivation of the letters: it changes the one-body kernels and leaves the two-body kernels unchanged. For the field countings $\phi^{(0)}=\sum_fy_f\int(f,f^*)$ the change of $Q_1$ multiplies each arrow $f\to f'$ of the classical transformation by $y_{f'}-y_f$; for $\phi^{(0)}=t\,S^{(1)}(\Delta)$ it is a multiple of the classical transformation proportional to $\lambda_x$.
\item (ii) The coupling direction leaves $Q_1$, $S_1$, $K_1$ unchanged. The outer direction $G_B$ changes $Q_1$ along $S_+\otimes W$ by $-a$ and along $S_-\otimes W^*$ by $+a$ times the classical transformation, $S_1$ along $S_+\otimes W^*$ by $+a$ and along $S_-\otimes W$ by $-a$ times it, and leaves $K_1$, $D_1$, $P_1$, $M_1$, $R_1$ unchanged.
\item (iii) The free directions of $\hbar F_1^{(1)}(x)$ change the one-body kernels: along $\Sigma$ by Proposition \ref{prop:c11-freexi} (iii); along $\Delta$ and $P_b$ through the sector coefficients of Definition \ref{def:c11-free}, which enter $D_1$, $P_1$ and, through $P_1$, the relations of Proposition \ref{prop:c11-S} and Corollary \ref{cor:c11-K}; along $K_b$ through the $s$-closed directions of degree $-1$ quadratic in the Killing data. $\hbar F_1^{(0)}$ enters the one-body kernels through \eqref{eq:c11-arity1xi} and through the trees with $S^{(1)}(x)$ (Remark \ref{rem:c10-scheme1} (a), (b)).
\item (iv) A finite part $\psi'$ of the operators changes the component along $x$ by $[\{S^{(1)}(x),-\},\psi']$, on $O$ by $\{S^{(1)}(x),\psi'(O)\}-\psi'(\{S^{(1)}(x),O\})$. For $\psi'=\tfrac a2\{S^{(1)}(\Delta),-\}$ this is $\tfrac a2\lambda_x\{S^{(1)}(x),-\}$: $+a$, $-\tfrac a2$, $+\tfrac a2$ times the classical transformation along $K_b$, $\Sigma$, $\Sigma^*$ (Proposition \ref{prop:c9-scheme}), a change of the one-body kernels. A finite part $\psi'$ that vanishes on $O$ changes the component of $O$ along $x$ by $-\psi'(\{S^{(1)}(x),O\})$: the finite parts of the operators of the Weyl weight of $\{S^{(1)}(x),O\}$ enter the two-body kernels, as the finite part of Lemma \ref{lem:c11-cocycle} on the trace bilinear and the shift of the Konishi entry in Remark \ref{rem:c11-twist}.
\item (v) The cutoff shape enters the finite parts of the defect cycles (Proposition \ref{prop:c10-shape}, Lemma \ref{lem:c10-special}), of the special-edge cycles of $L_1^{(0)}$ and of the finite special vertices; since the counterterms are unique up to finite local terms [C07, Thm B], a change of the shape is a change of $\hbar F_1$ and of $\psi$, a combination of (i)--(iv).
\item (vi) Under the hypotheses of Theorem \ref{thm:c11-arity2}, in the translation-covariant scheme at $\theta=\theta_\lambda=\theta_V=-\tfrac13$, $P_1\cong0$ on the families of Definition \ref{def:c11-families} (Corollary \ref{cor:c11-D1fam}), and the condition $P_1=0$ of Definition \ref{def:c8-radial} holds there with $\psi'=0$: the components along $K_b$, $\Sigma$ and $\Sigma^*$ are those of the scheme. For equal coefficients $\theta_P\ne\theta$ along $P_b$, the finite part with $a=(\theta-\theta_P)$ units gives $P_1\cong0$ and changes the components along $K_b$, $\Sigma$, $\Sigma^*$ by $a$, $-\tfrac a2$, $\tfrac a2$ times the classical transformation (Corollary \ref{cor:c11-theta} (iv)).
\item (vii) Under (i), (iv) and the coupling direction the representation of $\mathfrak{psl}(4|4)$ on the cohomology of $\Lambda^{(0)}$ changes by an isomorphism, under $G_B$ by the automorphism $e^{-a\hbar B}$, and the spectrum of $D_1$ in a fixed classical Weyl weight is invariant (Theorem \ref{thm:c8-invariants}); the class $[Q^\xi_1]$ on the cohomology of $Q^\xi_0$, its kernel, image and rank are invariant under (i), (ii) and (iv) (Theorem \ref{thm:c8-twist}).
\end{itemize}
\end{proposition}

\begin{proof}
(i), (ii) Proposition \ref{prop:c8-laws} (a), (c); the bracket with a local functional is a derivation in its second argument, so it acts letter by letter; for $\phi^{(0)}=t\,S^{(1)}(\Delta)$, $\{S^{(1)}(x),S^{(1)}(\Delta)\}$ is $S^{(1)}$ of $\pm[x,\Delta]=\mp\lambda_xx$ by Lemma \ref{lem:c2-components}, with $S^{(2)}(x,\Delta)=0$ (Definition \ref{def:c2-S2}). (iii) Proposition \ref{prop:c11-freexi}, Definition \ref{def:c11-free}, Remark \ref{rem:c10-scheme1}. (iv) Proposition \ref{prop:c8-laws} (b) and Proposition \ref{prop:c9-scheme}, with $[\{S^{(1)}(x),-\},\psi'](O)=\{S^{(1)}(x),\psi'(O)\}-\psi'(\{S^{(1)}(x),O\})$. (v) Remark \ref{rem:c10-scheme1} (d); the finite counterterm and the finite parts of the operators absorb the change. (vi) Corollaries \ref{cor:c11-D1fam} and \ref{cor:c11-theta} (iv). (vii) Theorems \ref{thm:c8-invariants} and \ref{thm:c8-twist}.
\end{proof}

\begin{lemma}[The trace bilinear in a cocycle scheme]
\label{lem:c11-cocycle}
Let $O=\int\omega\,(\iota_{w_1}A,\iota_{w_2}A)$, $\xi\in S_+\otimes W$ or $\xi\in S_-\otimes W^*$, and let $\mathcal B_1,\dots,\mathcal B_4$ be the four functionals of Corollary \ref{cor:c7-Q},
\[
\begin{aligned}
&\mathcal B_1:=\{S^{(1)}(\xi),O\},\qquad \mathcal B_2:=(w_1,w_2)\int\omega\,\big(\rho(A|_{W_{\mathbb R}^\perp})\xi,\lambda\big),\\
&\mathcal B_3:=(w_1,w_2)\int\omega\,\big(\rho(A|_V)\xi,\lambda\big),\qquad \mathcal B_4:=(w_1,w_2)\int\omega\,\big(c,(\xi,\lambda^*)\big),
\end{aligned}
\]
$\mathcal B_1$ being the first functional listed there (Lemma \ref{lem:c11-lines} (i)). Write the two-body part of $Q_1(O)$ as $a_1\mathcal B_1+a_2\mathcal B_2+a_3\mathcal B_3+a_4\mathcal B_4$.
\begin{itemize}
\item (i) If $L_1^{(0)}=0$ on $O$ and on $\mathcal B_1$, then $a_3=a_4=0$, and the one-body terms of $Q_1(O)$ are multiples of $\mathcal B_1$.
\item (ii) If $L_1^{(0)}(O)=2\kappa_O(w_1,w_2)\int\omega\,(dc,A|_V)$ and the part of $L_1^{(0)}(\mathcal B_1)$ with the two letters $c$, $\lambda$ is $-\kappa_{\mathcal B}\,s\mathcal B_3$, the finite part $\psi(O)=\kappa_O(w_1,w_2)\int\omega\,(A|_V,A|_V)$, $\psi(\mathcal B_1)=\kappa_{\mathcal B}\mathcal B_3$ removes both and changes $Q_1(O)$ by $(2\kappa_O-\kappa_{\mathcal B})\mathcal B_3$. At $\psi=0$ the part with the letters $c$, $\lambda$ of (R1) of Proposition \ref{prop:c11-relations} is the identity $a_3+2\kappa_O-\kappa_{\mathcal B}=0$.
\end{itemize}
\end{lemma}

\begin{proof}
(i) By (R1), $s\,Q_1(O)=0$. $s\mathcal B_1=0$, since $[s,\{S^{(1)}(\xi),-\}]=\{\{S,S^{(1)}(\xi)\},-\}=0$ at tree level and $sO=0$; $s\mathcal B_2=0$, since $Q$ annihilates the coordinate functionals of $A|_{W_{\mathbb R}^\perp}$ and $\lambda$ and the cubic parts of $s$ on them are the gauge transformation with parameter $c$, which leaves the invariant pairing fixed. In the sector with the two letters $c$, $\lambda$, $s\mathcal B_3=-(w_1,w_2)\int\omega\,(\rho(dc)\xi,\lambda)$ and $s\mathcal B_4=(w_1,w_2)\int\omega\,(c,(\xi,{\not d}\lambda))$, whose densities carry the derivative on different letters and are linearly independent for arbitrary $\omega$; hence $a_3=a_4=0$. The one-body terms are $-\{\hbar F_1^{(1)}(\xi),O\}$; a term of $\hbar F_1^{(1)}(\xi)$ acting on the letters $A|_{W_{\mathbb R}^\perp}$ of $O$ has the leg $A^*|_{W_{\mathbb R}^\perp}$, and by Proposition \ref{prop:c11-freexi} (i) it is a multiple of $T_2$, which gives a multiple of $\mathcal B_1$. (ii) $Q\int\omega\,(A|_V,A|_V)=-2\int\omega\,(dc,A|_V)$ (Remark \ref{rem:c5-cocycle}), and the cubic part of $s\int\omega\,(A|_V,A|_V)$ vanishes by the invariance of $(-,-)$; $L_1$ changes by $[s,\psi]$ (Lemma \ref{lem:c10-psi}), the component along $\xi$ by $[\{S^{(1)}(\xi),-\},\psi]$ (Proposition \ref{prop:c8-laws} (b)), and $\{S^{(1)}(\xi),\int\omega\,(A|_V,A|_V)\}=2\int\omega\,(\Gamma(\xi,\lambda)|_V,A|_V)$ with $(\Gamma(\xi,\lambda)|_V,a)=(\rho(a)\xi,\lambda)$ [ESW, (28)]. In the shifted scheme (i) gives $a_3+2\kappa_O-\kappa_{\mathcal B}=0$; at $\psi=0$ the part with the letters $c$, $\lambda$ of (R1), $s\,Q_1(O)+L_1^{(0)}(\mathcal B_1)+\{S^{(1)}(\xi),L_1^{(0)}(O)\}=0$, is the same equation, since there $s\,Q_1(O)=a_3\,s\mathcal B_3$, $L_1^{(0)}(\mathcal B_1)=-\kappa_{\mathcal B}\,s\mathcal B_3$ and $\{S^{(1)}(\xi),L_1^{(0)}(O)\}=2\kappa_O\,s\mathcal B_3$.
\end{proof}

\begin{theorem}[$Q_1$ on the scalar bilinears]
\label{thm:c11-bilinears}
Let $O=\int\omega\,(\iota_{w_1}A,\iota_{w_2}A)$, $\xi\in S_+\otimes W$ or $\xi\in S_-\otimes W^*$, and $\mathcal B_1,\dots,\mathcal B_4$ as in Lemma \ref{lem:c11-cocycle}. At $\psi=0$:
\begin{itemize}
\item (i) the two-body part of $Q_1(O)$ is
\begin{equation}
-2L\,\mathcal B_1+4L\,\mathcal B_2+4L\,\mathcal B_3+0\cdot\mathcal B_4\ \text{units},\qquad L=\log\tfrac32,
\label{eq:c11-Q1O}
\end{equation}
without poles and the same for both chiralities; its trace part is $4L$ units times $(w_1,w_2)\int\omega\,(\rho(A)\xi,\lambda)$, with the whole connection $A=A|_V+A|_{W_{\mathbb R}^\perp}$;
\item (ii) on the traceless part, $(w_1,w_2)=0$,
\[
Q_1(O)=c_Q\,\{S^{(1)}(\xi),O\},\qquad c_Q=-2\log\tfrac32-t_\perp\ \text{units},
\]
with $t_\perp$ units the coefficient of $T_2$ in $\hbar F_1^{(1)}(\xi)$, which \eqref{eq:c11-arity1xi} fixes up to the free direction $\theta^\xi_\perp$ of Proposition \ref{prop:c11-freexi}; $c_Q$ moves by $-\theta^\xi_\perp$ under that direction, by $\mp a$ under $G_B$ along $S_+\otimes W$ and $S_-\otimes W^*$, by $-\tfrac a2$ under $\psi'=\tfrac a2\{S^{(1)}(\Delta),-\}$, and under the field countings as in Proposition \ref{prop:c11-freedom} (i);
\item (iii) $L_1^{(0)}(O)=-(2+4L)$ units times $(w_1,w_2)\int\omega\,(dc,A|_V)$, and the part of $L_1^{(0)}(\mathcal B_1)$ with the letters $c$, $\lambda$ is $2$ units times $s\mathcal B_3$; Lemma \ref{lem:c11-cocycle} (ii) applies with $\kappa_O=-(1+2L)$ units and $\kappa_{\mathcal B}=-2$ units, and in the cocycle scheme
\[
\psi(O)=-(1+2L)\,\text{units}\cdot(w_1,w_2)\int\omega\,(A|_V,A|_V),\qquad \psi(\mathcal B_1)=-2\,\text{units}\cdot\mathcal B_3,
\]
the two-body part of $Q_1(O)$ is $-2L\,\mathcal B_1+4L\,\mathcal B_2$ units, $a_3=4L+2\kappa_O-\kappa_{\mathcal B}=0$;
\item (iv) at $\psi=0$, for $\xi\in S_+\otimes W$, the part with the letters $c$, $\lambda|_{S_-\otimes W^*}$ of (R1) is the identity
\[
\begin{aligned}
&s\,Q_1(O)+L_1^{(0)}\big(\{S^{(1)}(\xi),O\}\big)+\big\{S^{(1)}(\xi),L_1^{(0)}(O)\big\}\\
&\qquad=\big(-4L-2+(2+4L)\big)\,\text{units}\cdot(w_1,w_2)\int\omega\,\big(\rho(dc)\xi,\lambda\big)=0 .
\end{aligned}
\]
\end{itemize}
\end{theorem}

\begin{proof}
(i) By Theorem \ref{thm:c11-classes} (iii) the two-body part is $\Upsilon_\xi$ on the single pair $(A|_{W_{\mathbb R}^\perp},A|_{W_{\mathbb R}^\perp})$ of $O$, whose leg sets are those of Table \ref{tab:c11-defectpairs}: $\lambda|_{S_-\otimes W^*}$, $A|_V\,\lambda|_{S_-\otimes W^*}$, $A|_{W_{\mathbb R}^\perp}\lambda|_{S_+\otimes W}$ and $c\,\lambda^*|_{S_+\otimes W^*}$ for $\xi\in S_+\otimes W$. The single output $Y$ has the colour factor $(1\otimes\mathrm{ad}_Y)(-,-)^{-1}$, which the colour pairing of $O$ contracts to $\operatorname{tr}_{\mathfrak g}\mathrm{ad}_Y=0$. The other three leg sets have $D=0$, so they carry no poles and no finite special vertex (Proposition \ref{prop:c11-oddbound}), and they give $\mathcal B_1$, $\mathcal B_2$ (outputs $A|_{W_{\mathbb R}^\perp}\lambda$), $\mathcal B_3$ (outputs $A|_V\lambda$) and $\mathcal B_4$ (outputs $c\,\lambda^*$). Their chains have three edges and, at $D=0$, constant integrands, so each chain gives a rational multiple of $\tfrac12\log\tfrac32$ (Theorem \ref{thm:c11-oddsimplex} (iv)); the exact evaluation of the chains in all placements, in both ways of Remark \ref{rem:c11-oddmethod}, gives \eqref{eq:c11-Q1O}, with the same coefficient for $\mathcal B_2$ and $\mathcal B_3$. (ii) By Corollary \ref{cor:c7-Q} the functionals at $(w_1,w_2)=0$ are the multiples of $\mathcal B_1$; the one-body term is $-t_\perp\mathcal B_1$ units (proof of Lemma \ref{lem:c11-cocycle} (i)); the changes are Propositions \ref{prop:c11-freexi} and \ref{prop:c11-freedom}. (iii) $L_1^{(0)}$ is given by the arity-zero two-body kernels, the cycles with one special edge of the proof of Theorem \ref{thm:c10-brst}, of the pairs $(A|_{W_{\mathbb R}^\perp},A|_{W_{\mathbb R}^\perp})$ of $O$ and $(A|_{W_{\mathbb R}^\perp},\lambda)$ of $\mathcal B_1$, evaluated in the same way; then Lemma \ref{lem:c11-cocycle} (ii). (iv) With $s\mathcal B_3=-(w_1,w_2)\int\omega\,(\rho(dc)\xi,\lambda)$ in this sector, the three terms are $4L\,s\mathcal B_3$, $2\,s\mathcal B_3$ and $2\kappa_O\,s\mathcal B_3=-(2+4L)\,s\mathcal B_3$ units (proof of Lemma \ref{lem:c11-cocycle} (ii)).
\end{proof}

\begin{proposition}[$S_1$ and $K_1$ on the scalar bilinears]
\label{prop:c11-SKbil}
Let $O=\int\omega\,(\iota_{w_1}A,\iota_{w_2}A)$ and $\psi=0$.
\begin{itemize}
\item (i) $\Upsilon^0_\eta(A|_{W_{\mathbb R}^\perp},A|_{W_{\mathbb R}^\perp})$ has the single leg set $\lambda$, which vanishes on the colour singlet, and $s_0$ has zero bracket with a letter $A|_{W_{\mathbb R}^\perp}$. Hence the two-body part of $S_1(O)$ at the point $x$ of the density is the two-body part of $Q_1(O)$ along the Killing spinor $\rho(x)\eta$, given by \eqref{eq:c11-Q1O} along $S_-\otimes W^*$ for $\eta\in S_+\otimes W^*$ and along $S_+\otimes W$ for $\eta\in S_-\otimes W$; the term linear in ${\not d}(\rho(x)\eta)=4\eta$ is absent (Corollary \ref{cor:c7-S}).
\item (ii) $\Upsilon^0_{K_b}(A|_{W_{\mathbb R}^\perp},A|_{W_{\mathbb R}^\perp})$ has the leg sets $\emptyset$, of superficial degree $1$, and $A|_V$, both zero on the colour singlet. In a scheme in which $\hbar F_1^{(1)}(K_b)$ decomposes as $S^{(1)}(K_b)$ in Theorem \ref{thm:c11-oddsimplex} (iii), on the part with two legs $A|_{W_{\mathbb R}^\perp}$,
\[
K_1(O)=v(x)^\mu\,P_1(O)\big|_{P_{e_\mu}}-(b,x)\,D^0_1(O),
\]
with $D^0_1$ the translation-invariant part of $D_1$ (Definition \ref{def:c11-canonical}); the rotation part and the intrinsic part vanish. With Theorems \ref{thm:c9-D1} and \ref{thm:c9-P1} at $\theta_P=\theta$,
\[
\begin{aligned}
&K_1(O)=-\big(\tfrac13+\theta\big)\,\text{units}\cdot\{S^{(1)}(K_b),O\}+\begin{cases}0&\text{(traceless part)},\\ 48\,\text{units}\cdot(b,x)\,O&\text{(trace)},\end{cases}\\
&\{S^{(1)}(K_b),O\}=v^\mu\partial_\mu O-4(b,x)\,O,
\end{aligned}
\]
so that at $\theta=-\tfrac13$, $K_1(O)=-(b,x)D_1(O)$: $0$ on the traceless bilinears and $48$ units times $(b,x)\,O$ on the trace, the relation of Proposition \ref{prop:c8-primary} with $\kappa=-24$ units.
\end{itemize}
\end{proposition}

\begin{proof}
(i) Table \ref{tab:c11-defectcounts} and the evaluation of the leg set: a single output $Y$ gives $\operatorname{tr}_{\mathfrak g}\mathrm{ad}_Y=0$ on the colour singlet; $s_0$ has the legs $A|_{W_{\mathbb R}^\perp}$, $\lambda^*$, and its bracket with a letter needs the leg $A^*|_{W_{\mathbb R}^\perp}$. Theorem \ref{thm:c11-oddsimplex} (ii) gives the position part. (ii) The leg set $A|_V$ has a single output; the empty leg set is homogeneous in $\varepsilon$ of degree $-\tfrac12$ (Lemma \ref{lem:c10-self}). Theorem \ref{thm:c11-oddsimplex} (iii) for the two-body parts, with the one-body term decomposed in the same way; by Corollary \ref{cor:c7-K} the term in the second derivatives of $v$ and the term in the antisymmetric part of $dv$ vanish on a scalar bilinear. By \eqref{eq:c9-P1}, $P_1(O)|_{P_{e_\mu}}=-(\tfrac13+\theta_P)$ units times $\partial_\mu O$; by \eqref{eq:c9-D1}, with $\{S^{(1)}(\Delta),O\}=2x^\mu\partial_\mu O+4O$, $D^0_1(O)$ at the point of the density is $-4(\tfrac13+\theta)$ units times $O$, plus $-48$ units times $O$ on the trace; and $\{S^{(1)}(K_b),O\}=v^\mu\partial_\mu O+\tfrac24(\operatorname{div}v)O$ (Definition \ref{def:c7-vdot}). Proposition \ref{prop:c8-primary}: $K_1(O)=-2(b,x)\kappa\,O$ for $D_1(O)=2\kappa\,O$.
\end{proof}

\begin{corollary}[The dilatation and the translations by the defect cycles]
\label{cor:c11-DPdefect}
On every local functional $O$ with $sO=0$, with the Koszul signs of \eqref{eq:c10-decomp},
\[
\begin{aligned}
&\sum_{a<b}\pm\Upsilon_\Delta(X_a,X_b)=\sum_{a<b}\pm\big(4\,\mathrm{lc}(X_aX_b)+\text{composite pair kernel along }\Delta\big),\\
&\sum_{a<b}\pm\Upsilon_{P_b}(X_a,X_b)=\sum_{a<b}\pm\big(\text{composite pair kernel along }P_b\big),
\end{aligned}
\]
and the one-body terms of the two forms agree. On the part with two legs $A|_{W_{\mathbb R}^\perp}$ of the scalar bilinears, at the point of the density, $\Upsilon^0_\Delta$ is $\tfrac{40}3$ units times $O$ on the traceless part and $-\tfrac{104}3$ units times $O$ on the trace, and $\Upsilon_{P_b}$ is $-\tfrac13$ units times $\partial_bO$; with the position part $2x^\mu\Upsilon_{P_{e_\mu}}$ this is $\tfrac{44}3O-\tfrac13\{S^{(1)}(\Delta),O\}$ and $-\tfrac{100}3O-\tfrac13\{S^{(1)}(\Delta),O\}$ units, the first row of the table of Theorem \ref{thm:c9-D1}, and the value $-\tfrac13$ of Theorem \ref{thm:c9-P1}.
\end{corollary}

\begin{proof}
Both forms compute the unique local representative (Theorem \ref{thm:c8-L} (i)): for $\Delta$ this is Proposition \ref{prop:c11-routes}, whose route of the special edges is \eqref{eq:c11-decompx}, and for $P_b$ the same argument applies with Corollary \ref{cor:c11-D1local} and the one-body term $-\{\hbar F_1^{(1)}(P_b),-\}$ in both forms. On the bilinears, the defect cycles with $S^{(1)}(\Delta)$ and $S^{(1)}(P_b)$ in all their placements, evaluated by Theorem \ref{thm:c11-oddsimplex}, give the stated values, with $\{S^{(1)}(\Delta),O\}=2x^\mu\partial_\mu O+4O$ at the point of the density: $\tfrac{40}3+\tfrac43=\tfrac{44}3$ and $-\tfrac{104}3+\tfrac43=-\tfrac{100}3$.
\end{proof}

